\section{Smooth morphisms}
\label{section-smooth}

\noindent
Let $f : X \to Y$ be a continuous map of topological spaces.
Consider the following condition:
For every $x \in X$ there exist open neighbourhoods
$x \in U \subset X$ and $f(x) \in V \subset Y$, and an integer $d$
such that $f(U) \subset V$ and such that we obtain
a commutative diagram
$$
\xymatrix{
X \ar[d] &
U \ar[l] \ar[d] \ar[r]_-\pi &
V \times \mathbf{R}^d \ar[ld] \\
Y & V \ar[l]
}
$$
where $\pi$ is a homeomorphism onto an open subset.
Smooth morphisms of schemes are the analogue of these maps
in the category of schemes. See
Lemma \ref{lemma-smooth-locally-standard-smooth}
and
Lemma \ref{lemma-smooth-etale-over-affine-space}.

\medskip\noindent
Contrary to expectations (perhaps) the notion
of a smooth ring map is not defined solely in terms
of the module of differentials. Namely, recall that
$R \to A$ is a {\it smooth ring map} if $A$ is of finite presentation over $R$
and if the naive cotangent complex of $A$ over $R$ is quasi-isomorphic
to a projective module placed in degree $0$, see
Algebra, Definition \ref{algebra-definition-smooth}.

\begin{definition}
\label{definition-smooth}
Let $f : X \to S$ be a morphism of schemes.
\begin{enumerate}
\item We say that $f$ is {\it smooth at $x \in X$} if
there exist an affine open neighbourhood $\Spec(A) = U \subset X$
of $x$ and affine open $\Spec(R) = V \subset S$
with $f(U) \subset V$ such that the induced ring map
$R \to A$ is smooth.
\item We say that $f$ is {\it smooth} if it is smooth at every point of $X$.
\item A morphism of affine schemes $f : X \to S$
is called {\it standard smooth} if there exists a standard smooth ring
map $R \to R[x_1, \ldots, x_n]/(f_1, \ldots, f_c)$ (see
Algebra, Definition \ref{algebra-definition-standard-smooth})
such that $X \to S$ is isomorphic to
$$
\Spec(R[x_1, \ldots, x_n]/(f_1, \ldots, f_c)) \to \Spec(R).
$$
\end{enumerate}
\end{definition}

\noindent
A pleasing feature of this definition is that the set of points
where a morphism is smooth is automatically open.

\medskip\noindent
Note that there is no separation or quasi-compactness hypotheses in the
definition. Hence the question of being smooth is local in nature on
the source. Here is the precise result.

\begin{lemma}
\label{lemma-smooth-characterize}
Let $f : X \to S$ be a morphism of schemes.
The following are equivalent
\begin{enumerate}
\item The morphism $f$ is smooth.
\item For every affine opens $U \subset X$, $V \subset S$
with $f(U) \subset V$ the ring map
$\mathcal{O}_S(V) \to \mathcal{O}_X(U)$ is smooth.
\item There exists an open covering $S = \bigcup_{j \in J} V_j$
and open coverings $f^{-1}(V_j) = \bigcup_{i \in I_j} U_i$ such
that each of the morphisms $U_i \to V_j$, $j\in J, i\in I_j$
is smooth.
\item There exists an affine open covering $S = \bigcup_{j \in J} V_j$
and affine open coverings $f^{-1}(V_j) = \bigcup_{i \in I_j} U_i$ such
that the ring map $\mathcal{O}_S(V_j) \to \mathcal{O}_X(U_i)$ is
smooth, for all $j\in J, i\in I_j$.
\end{enumerate}
Moreover, if $f$ is smooth then for
any open subschemes $U \subset X$, $V \subset S$ with $f(U) \subset V$
the restriction $f|_U : U \to V$ is smooth.
\end{lemma}

\begin{proof}
This follows from Lemma \ref{lemma-locally-P} if we show that
the property ``$R \to A$ is smooth'' is local.
We check conditions (a), (b) and (c) of Definition
\ref{definition-property-local}.
By Algebra, Lemma \ref{algebra-lemma-base-change-smooth}
being smooth is stable under base change and hence
we conclude (a) holds. By
Algebra, Lemma \ref{algebra-lemma-compose-smooth}
being smooth is stable under composition and for any ring
$R$ the ring map $R \to R_f$ is (standard) smooth.
We conclude (b) holds. Finally, property (c) is true
according to Algebra, Lemma \ref{algebra-lemma-locally-smooth}.
\end{proof}

\noindent
The following lemma characterizes a smooth morphism as a
flat, finitely presented morphism with smooth fibres.
Note that schemes smooth over a field are discussed in more detail in
Varieties, Section \ref{varieties-section-smooth}.

\begin{lemma}
\label{lemma-smooth-flat-smooth-fibres}
Let $f : X \to S$ be a morphism of schemes.
If $f$ is flat, locally of finite presentation, and all
fibres $X_s$ are smooth, then $f$
is smooth.
\end{lemma}

\begin{proof}
Follows from Algebra, Lemma \ref{algebra-lemma-flat-fibre-smooth}.
\end{proof}

\begin{lemma}
\label{lemma-composition-smooth}
The composition of two morphisms which are smooth is smooth.
\end{lemma}

\begin{proof}
In the proof of Lemma \ref{lemma-smooth-characterize}
we saw that being smooth is a local property of ring maps.
Hence the first statement of the lemma follows from
Lemma \ref{lemma-composition-type-P} combined
with the fact that being smooth is a property of ring maps that is
stable under composition, see
Algebra, Lemma \ref{algebra-lemma-compose-smooth}.
\end{proof}

\begin{lemma}
\label{lemma-base-change-smooth}
The base change of a morphism which is smooth is smooth.
\end{lemma}

\begin{proof}
In the proof of Lemma \ref{lemma-smooth-characterize}
we saw that being smooth is a local property of ring maps.
Hence the lemma follows from
Lemma \ref{lemma-composition-type-P} combined
with the fact that being smooth is a property of ring maps that is
stable under base change, see
Algebra, Lemma \ref{algebra-lemma-base-change-smooth}.
\end{proof}

\begin{lemma}
\label{lemma-open-immersion-smooth}
Any open immersion is smooth.
\end{lemma}

\begin{proof}
This is true because an open immersion is a local isomorphism.
\end{proof}

\begin{lemma}
\label{lemma-smooth-syntomic}
A smooth morphism is syntomic.
\end{lemma}

\begin{proof}
See Algebra, Lemma \ref{algebra-lemma-smooth-syntomic}.
\end{proof}

\begin{lemma}
\label{lemma-smooth-locally-finite-presentation}
A smooth morphism is locally of finite presentation.
\end{lemma}

\begin{proof}
True because a smooth ring map is of finite presentation by
definition.
\end{proof}

\begin{lemma}
\label{lemma-smooth-flat}
A smooth morphism is flat.
\end{lemma}

\begin{proof}
Combine Lemmas \ref{lemma-syntomic-flat} and \ref{lemma-smooth-syntomic}.
\end{proof}

\begin{lemma}
\label{lemma-smooth-open}
A smooth morphism is universally open.
\end{lemma}

\begin{proof}
Combine
Lemmas \ref{lemma-smooth-flat},
\ref{lemma-smooth-locally-finite-presentation}, and
\ref{lemma-fppf-open}.
Or alternatively, combine
Lemmas \ref{lemma-smooth-syntomic},
\ref{lemma-syntomic-open}.
\end{proof}

\noindent
The following lemma says locally any smooth morphism is standard smooth.
Hence we can use standard smooth morphisms as a {\it local model}
for a smooth morphism.

\begin{lemma}
\label{lemma-smooth-locally-standard-smooth}
\begin{slogan}
Smooth morphisms are locally standard smooth.
\end{slogan}
Let $f : X  \to S$ be a morphism of schemes.
Let $x \in X$ be a point.
Let $V \subset S$ be an affine open neighbourhood of $f(x)$.
The following are equivalent
\begin{enumerate}
\item The morphism $f$ is smooth at $x$.
\item There exists an affine open $U \subset X$,
with $x \in U$ and $f(U) \subset V$ such that the
induced morphism $f|_U : U \to V$ is standard smooth.
\end{enumerate}
\end{lemma}

\begin{proof}
Follows from the definitions and
Algebra, Lemmas \ref{algebra-lemma-standard-smooth}
and \ref{algebra-lemma-smooth-syntomic}.
\end{proof}

\begin{lemma}
\label{lemma-smooth-omega-finite-locally-free}
Let $f : X \to S$ be a morphism of schemes.
Assume $f$ is smooth.
Then the module of differentials $\Omega_{X/S}$ of $X$ over $S$
is finite locally free and
$$
\text{rank}_x(\Omega_{X/S}) = \dim_x(X_{f(x)})
$$
for every $x \in X$.
\end{lemma}

\begin{proof}
The statement is local on $X$ and $S$.
By Lemma \ref{lemma-smooth-locally-standard-smooth}
above we may assume that $f$ is a standard smooth morphism of affines.
In this case the result follows from
Algebra, Lemma \ref{algebra-lemma-standard-smooth}
(and the definition of a relative global complete intersection, see
Algebra,
Definition \ref{algebra-definition-relative-global-complete-intersection}).
\end{proof}

\noindent
Lemma \ref{lemma-smooth-omega-finite-locally-free}
says that the following definition makes sense.

\begin{definition}
\label{definition-smooth-relative-dimension}
Let $d \geq 0$ be an integer. We say a morphism of schemes $f : X \to S$
is {\it smooth of relative dimension $d$} if $f$ is smooth and
$\Omega_{X/S}$ is finite locally free of constant rank $d$.
\end{definition}

\noindent
In other words, $f$ is smooth and the nonempty fibres are equidimensional
of dimension $d$. By Lemma \ref{lemma-smooth-at-point} below this is also
the same as requiring: (a) $f$ is locally of finite presentation, (b) $f$ is
flat, (c) all nonempty fibres equidimensional of dimension $d$, and (d)
$\Omega_{X/S}$ finite locally free of rank $d$. It is not enough to simply
assume that $f$ is flat, of finite presentation, and $\Omega_{X/S}$ is
finite locally free of rank $d$. A counter example is given by
$\Spec(\mathbf{F}_p[t]) \to \Spec(\mathbf{F}_p[t^p])$.

\medskip\noindent
Here is a differential criterion of smoothness at a point.
There are many variants of this result
all of which may be useful at some point. We will just add them
here as needed.

\begin{lemma}
\label{lemma-smooth-at-point}
Let $f : X \to S$ be a morphism of schemes.
Let $x \in X$.
Set $s = f(x)$.
Assume $f$ is locally of finite presentation.
The following are equivalent:
\begin{enumerate}
\item The morphism $f$ is smooth at $x$.
\item The local ring map $\mathcal{O}_{S, s} \to \mathcal{O}_{X, x}$
is flat and $X_s \to \Spec(\kappa(s))$ is smooth at $x$.
\item The local ring map $\mathcal{O}_{S, s} \to \mathcal{O}_{X, x}$
is flat and the $\mathcal{O}_{X, x}$-module $\Omega_{X/S, x}$
can be generated by at most $\dim_x(X_{f(x)})$ elements.
\item The local ring map $\mathcal{O}_{S, s} \to \mathcal{O}_{X, x}$
is flat and the $\kappa(x)$-vector space
$$
\Omega_{X_s/s, x} \otimes_{\mathcal{O}_{X_s, x}} \kappa(x) =
\Omega_{X/S, x} \otimes_{\mathcal{O}_{X, x}} \kappa(x)
$$
can be generated by at most $\dim_x(X_{f(x)})$ elements.
\item There exist affine opens $U \subset X$,
and $V \subset S$ such that $x \in U$, $f(U) \subset V$ and the
induced morphism $f|_U : U \to V$ is standard smooth.
\item There exist affine opens $\Spec(A) = U \subset X$
and $\Spec(R) = V \subset S$ with $x \in U$ corresponding
to $\mathfrak q \subset A$, and $f(U) \subset V$
such that there exists a presentation
$$
A = R[x_1, \ldots, x_n]/(f_1, \ldots, f_c)
$$
with
$$
g =
\det
\left(
\begin{matrix}
\partial f_1/\partial x_1 &
\partial f_2/\partial x_1 &
\ldots &
\partial f_c/\partial x_1 \\
\partial f_1/\partial x_2 &
\partial f_2/\partial x_2 &
\ldots &
\partial f_c/\partial x_2 \\
\ldots & \ldots & \ldots & \ldots \\
\partial f_1/\partial x_c &
\partial f_2/\partial x_c &
\ldots &
\partial f_c/\partial x_c
\end{matrix}
\right)
$$
mapping to an element of $A$ not in $\mathfrak q$.
\end{enumerate}
\end{lemma}

\begin{proof}
Note that if $f$ is smooth at $x$, then we see from Lemma
\ref{lemma-smooth-locally-standard-smooth} that (5) holds, and (6) is a slightly
weakened version of (5). Moreover, $f$ smooth implies that the ring
map $\mathcal{O}_{S, s} \to \mathcal{O}_{X, x}$ is flat (see
Lemma \ref{lemma-smooth-flat}) and that $\Omega_{X/S}$ is
finite locally free of rank equal to
$\dim_x(X_s)$ (see Lemma \ref{lemma-smooth-omega-finite-locally-free}).
Thus (1) implies (3) and (4). By Lemma \ref{lemma-base-change-smooth}
we also see that (1) implies (2).

\medskip\noindent
By Lemma \ref{lemma-base-change-differentials}
the module of differentials $\Omega_{X_s/s}$ of the fibre $X_s$
over $\kappa(s)$ is the pullback of the module of differentials
$\Omega_{X/S}$ of $X$ over $S$. Hence the displayed equality in
part (4) of the lemma. By Lemma \ref{lemma-finite-type-differentials}
these modules are of finite type. Hence the minimal number of
generators of the modules
$\Omega_{X/S, x}$ and $\Omega_{X_s/s, x}$ is the same and equal to the
dimension of this $\kappa(x)$-vector space by Nakayama's Lemma
(Algebra, Lemma \ref{algebra-lemma-NAK}). This in particular shows that
(3) and (4) are equivalent.

\medskip\noindent
Algebra, Lemma \ref{algebra-lemma-flat-fibre-smooth} shows that
(2) implies (1).
Algebra, Lemma \ref{algebra-lemma-characterize-smooth-over-field}
shows that (3) and (4) imply (2). Finally, (6) implies (5)
see for example Algebra, Example \ref{algebra-example-make-standard-smooth}
and (5) implies (1) by Algebra, Lemma \ref{algebra-lemma-standard-smooth}.
\end{proof}

\begin{lemma}
\label{lemma-set-points-where-fibres-smooth}
Let
$$
\xymatrix{
X' \ar[r]_{g'} \ar[d]_{f'} & X \ar[d]^f \\
S' \ar[r]^g & S
}
$$
be a cartesian diagram of schemes. Let $W \subset X$, resp.\ $W' \subset X'$
be the open subscheme of points where $f$, resp.\ $f'$ is smooth.
Then $W' = (g')^{-1}(W)$ if
\begin{enumerate}
\item $f$ is flat and locally of finite presentation, or
\item $f$ is locally of finite presentation and $g$ is flat.
\end{enumerate}
\end{lemma}

\begin{proof}
Assume first that $f$ locally of finite type. Consider the set
$$
T = \{x \in X \mid X_{f(x)}\text{ is smooth over }\kappa(f(x))\text{ at }x\}
$$
and the corresponding set $T' \subset X'$ for $f'$. Then we claim
$T' = (g')^{-1}(T)$. Namely, let $s' \in S'$ be a point, and let
$s = g(s')$. Then we have
$$
X'_{s'} =
\Spec(\kappa(s')) \times_{\Spec(\kappa(s))} X_s
$$
In other words the fibres of the base change are the base changes
of the fibres. Hence the claim is equivalent to
Algebra, Lemma \ref{algebra-lemma-smooth-field-change-local}.

\medskip\noindent
Thus case (1) follows because in case (1) $T$ is the (open) set of points
where $f$ is smooth by Lemma \ref{lemma-smooth-at-point}.

\medskip\noindent
In case (2) let $x' \in W'$. Then $g'$ is flat at $x'$
(Lemma \ref{lemma-base-change-module-flat}) and
$g \circ f'$ is flat at $x'$ (Lemma \ref{lemma-composition-module-flat}).
It follows that $f$ is flat at $x = g'(x')$
by Lemma \ref{lemma-flat-permanence}. On the other hand, since
$x' \in T'$ (Lemma \ref{lemma-base-change-smooth})
we see that $x \in T$. Hence $f$ is smooth at $x$ by
Lemma \ref{lemma-smooth-at-point}.
\end{proof}

\noindent
Here is a lemma that actually uses the vanishing of $H^{-1}$
of the naive cotangent complex for a smooth ring map.

\begin{lemma}
\label{lemma-triangle-differentials-smooth}
Let $f : X \to Y$, $g : Y \to S$ be morphisms of schemes.
Assume $f$ is smooth. Then
$$
0 \to f^*\Omega_{Y/S} \to \Omega_{X/S} \to \Omega_{X/Y} \to 0
$$
(see Lemma \ref{lemma-triangle-differentials}) is short exact.
\end{lemma}

\begin{proof}
The algebraic version of this lemma is the following:
Given ring maps $A \to B \to C$ with $B \to C$ smooth, then the sequence
$$
0 \to C \otimes_B \Omega_{B/A} \to \Omega_{C/A} \to \Omega_{C/B} \to 0
$$
of
Algebra, Lemma \ref{algebra-lemma-exact-sequence-differentials}
is exact. This is
Algebra, Lemma \ref{algebra-lemma-triangle-differentials-smooth}.
\end{proof}

\begin{lemma}
\label{lemma-differentials-relative-immersion-smooth}
Let $i : Z \to X$ be an immersion of schemes over $S$.
Assume that $Z$ is smooth over $S$. Then the
canonical exact sequence
$$
0 \to \mathcal{C}_{Z/X} \to i^*\Omega_{X/S} \to \Omega_{Z/S} \to 0
$$
of
Lemma \ref{lemma-differentials-relative-immersion}
is short exact.
\end{lemma}

\begin{proof}
The algebraic version of this lemma is the following:
Given ring maps $A \to B \to C$ with $A \to C$ smooth and $B \to C$
surjective with kernel $J$, then the sequence
$$
0 \to J/J^2 \to C \otimes_B \Omega_{B/A} \to \Omega_{C/A} \to 0
$$
of
Algebra, Lemma \ref{algebra-lemma-differential-seq}
is exact. This is
Algebra, Lemma \ref{algebra-lemma-differential-seq-smooth}.
\end{proof}

\begin{lemma}
\label{lemma-two-immersions-smooth}
Let
$$
\xymatrix{
Z \ar[r]_i \ar[rd]_j & X \ar[d] \\
& Y
}
$$
be a commutative diagram of schemes where $i$ and $j$ are immersions
and $X \to Y$ is smooth.
Then the canonical exact sequence
$$
0 \to  \mathcal{C}_{Z/Y} \to \mathcal{C}_{Z/X} \to i^*\Omega_{X/Y} \to 0
$$
of
Lemma \ref{lemma-two-immersions}
is exact.
\end{lemma}

\begin{proof}
The algebraic version of this lemma is the following:
Given ring maps $A \to B \to C$ with $A \to C$ surjective and $A \to B$
smooth, then the sequence
$$
0 \to I/I^2 \to J/J^2 \to C \otimes_B \Omega_{B/A} \to 0
$$
of
Algebra, Lemma \ref{algebra-lemma-application-NL}
is exact. This is
Algebra, Lemma \ref{algebra-lemma-application-NL-smooth}.
\end{proof}

\begin{lemma}
\label{lemma-smooth-permanence}
Let
$$
\xymatrix{
X \ar[rr]_f \ar[rd]_p & &
Y \ar[dl]^q \\
& S
}
$$
be a commutative diagram of morphisms of schemes. Assume that
\begin{enumerate}
\item $f$ is surjective, and smooth,
\item $p$ is smooth, and
\item $q$ is locally of finite presentation\footnote{In fact this
is implied by (1) and (2), see
Descent, Lemma \ref{descent-lemma-flat-finitely-presented-permanence}.
Moreover, it suffices to assume $f$ is surjective, flat and locally
of finite presentation, see
Descent, Lemma \ref{descent-lemma-smooth-permanence}.}.
\end{enumerate}
Then $q$ is smooth.
\end{lemma}

\begin{proof}
By Lemma \ref{lemma-flat-permanence} we see that $q$ is flat.
Pick a point $y \in Y$. Pick a point $x \in X$ mapping to $y$.
Suppose $f$ has relative dimension $a$ at $x$ and $p$ has relative
dimension $b$ at $x$. By Lemma \ref{lemma-smooth-omega-finite-locally-free}
this means that $\Omega_{X/S, x}$ is free of rank $b$ and $\Omega_{X/Y, x}$
is free of rank $a$. By the short exact sequence
of Lemma \ref{lemma-triangle-differentials-smooth}
this means that $(f^*\Omega_{Y/S})_x$ is free
of rank $b - a$. By Nakayama's Lemma this implies that
$\Omega_{Y/S, y}$ can be generated by $b - a$ elements.
Also, by Lemma \ref{lemma-dimension-fibre-at-a-point-additive} we see that
$\dim_y(Y_s) = b - a$. Hence we conclude that
$Y \to S$ is smooth at $y$ by Lemma \ref{lemma-smooth-at-point} part (2).
\end{proof}

\noindent
In the situation of the following lemma the image of $\sigma$ is
locally on $X$ cut out by a regular sequence, see
Divisors, Lemma \ref{divisors-lemma-section-smooth-regular-immersion}.

\begin{lemma}
\label{lemma-section-smooth-morphism}
Let $f : X \to S$ be a morphism of schemes.
Let $\sigma : S \to X$ be a section of $f$.
Let $s \in S$ be a point such that $f$ is smooth at $x = \sigma(s)$.
Then there exist affine open neighbourhoods
$\Spec(A) = U \subset S$ of $s$ and $\Spec(B) = V \subset X$
of $x$ such that
\begin{enumerate}
\item $f(V) \subset U$ and $\sigma(U) \subset V$,
\item with $I = \Ker(\sigma^\# : B \to A)$ the module $I/I^2$
is a free $A$-module, and
\item $B^\wedge \cong A[[x_1, \ldots, x_d]]$ as $A$-algebras where
$B^\wedge$ denotes the completion of $B$ with respect to $I$.
\end{enumerate}
\end{lemma}

\begin{proof}
Pick an affine open $U \subset S$ containing $s$
Pick an affine open $V \subset f^{-1}(U)$ containing $x$.
Pick an affine open $U' \subset \sigma^{-1}(V)$ containing $s$.
Note that $V' = f^{-1}(U') \cap V$ is affine as it is equal to the
fibre product $V' = U' \times_U V$. Then $U'$ and $V'$ satisfy (1).
Write $U' = \Spec(A')$ and $V' = \Spec(B')$. By
Algebra, Lemma \ref{algebra-lemma-section-smooth}
the module $I'/(I')^2$ is finite locally free as a $A'$-module.
Hence after replacing $U'$ by a smaller affine open $U'' \subset U'$
and $V'$ by $V'' = V' \cap f^{-1}(U'')$ we obtain the situation where
$I''/(I'')^2$ is free, i.e., (2) holds. In this case (3) holds also by
Algebra, Lemma \ref{algebra-lemma-section-smooth}.
\end{proof}

\noindent
The dimension of a scheme $X$ at a point $x$
(Properties, Definition \ref{properties-definition-dimension})
is just the dimension of $X$ at $x$ as a topological space, see
Topology, Definition \ref{topology-definition-Krull}.
This is not the dimension of the local ring $\mathcal{O}_{X,x}$, in general.

\begin{lemma}
\label{lemma-smoothness-dimension}
Let $f : X \to Y$ be a smooth morphism of locally Noetherian schemes.
For every point $x$ in $X$ with image $y$ in $Y$,
$$
\dim_x(X) = \dim_y(Y) + \dim_x(X_y),
$$
where $X_y$ denotes the fiber over $y$.
\end{lemma}

\begin{proof}
After replacing $X$ by an open neighborhood of $x$,
there is a natural number $d$ such that all fibers
of $X \to Y$ have dimension $d$ at every point, see
Lemma \ref{lemma-smooth-omega-finite-locally-free}.
Then $f$ is flat (Lemma \ref{lemma-smooth-flat}),
locally of finite type (Lemma \ref{lemma-smooth-locally-finite-presentation}),
and of relative dimension $d$. Hence the result follows from
Lemma \ref{lemma-rel-dimension-dimension}.
\end{proof}












